\chapter{从欧氏格点到光锥：LaMET 革命}

部分子分布是光锥关联；格点计算欧氏关联。三十年来两者之间的桥梁仅限
twist-2 算符的 Mellin 矩。2013 年 Xiangdong Ji 提出直接计算完整的
$x$ 依赖：把强子 boost 到大动量，测量等时空间关联，再经微扰匹配得到
光锥分布。本章收录定义该框架（LaMET）及其坐标空间孪生方法
（伪分布）的七篇论文。

% ----------------------------------------------------------------------
\section{奠基提案}
\papercard{Xiangdong Ji}{欧氏格点上的部分子物理}
{Phys.\ Rev.\ Lett.\ \textbf{110}, 262002 (2013)}
{arXiv:1305.1539 | DOI: 10.1103/PhysRevLett.110.262002}
\whyselected{本次挖掘的整个引文网络中最具影响力的单篇（理论语料约
$\sim$19/25；库内几乎每篇 PDF/DA/TMD 论文皆引用）。它本身就是本库
原文之一（\texttt{Parton Physics on Euclidean Lattice.pdf}），经核实
即为 PRL 正文。}
\physics{通过带 Wilson 线的等时、空间分离双局域算符定义\emph{准 PDF}：
\begin{equation}\label{eq:quasipdf}
  Q_\Gamma(z,P^z) = \langle P|\,\bar\psi(z\hat z)\,\Gamma\,
  W(z\hat z,0)\,\psi(0)\,|P\rangle ,
\end{equation}
它在格点上的计算与任何三点函数完全一样。大 $P^z$ 下光锥主导动力学地
涌现：单圈分析给出因子化
\begin{equation}\label{eq:l ametmatch}
  Q(x,P^z,\mu) = \int_x^\infty \frac{dy}{|x|}\;
  K\!\left(\frac{y}{x},\frac{\mu}{P^z}\right) q(y,\mu)
  + \mathcal{O}\!\left(\frac{\Lambda_{\rm QCD}^2}{(P^z)^2}\right),
\end{equation}
匹配核 $K$ 微扰可算且普适。Ji 在单圈检验了 $K$ 在短距离极限复现
Altarelli--Parisi 核，勾勒了到胶子、GPD、TMD 与 Wigner 分布的推广，
并指出功率修正就是该方法的代价。}
\impact{这是一个子领域的出生证明。本书其后每章要么实现、要么精化、
要么重整化、要么质疑公式~(\ref{eq:l ametmatch})。在本库内，准 PDF
数值计算（Lin 组）、伪 PDF 数值计算（HadStruc）、TMD 纲领（LPC）乃至
智能体工作流论文都可追溯到这一提案。}
\cardend

% ----------------------------------------------------------------------
\section{作为有效场论的 LaMET}
\papercard{Xiangdong Ji}{由大动量有效场论出发的部分子物理}
{Sci.\ China Phys.\ Mech.\ Astron.\ \textbf{57}, 1407--1412 (2014)}
{arXiv:1404.6680 | DOI: 10.1007/s11433-014-5450-3}
\whyselected{理论枢纽第二名（约 $\sim$15/99）：把大动量有效理论
（LaMET）表述为系统展开的正式文献，库内每篇方法论后期论文都引用它。}
\physics{高速态的行为如同受一个以 $\Lambda_{\rm QCD}/P^z$ 为展开参数
的有效理论支配：观测量分解为
\begin{equation}
  f(P^z) = \sum_n C_n(P^z,\mu)\otimes f_n(\mu)\,
  \left(\frac{\Lambda_{\rm QCD}}{P^z}\right)^{\!\omega_n},
\end{equation}
其中 $C_n$ 是普适匹配系数，$f_n$ 是一小族准算符的矩阵元。文章阐明使
该类"普适"的要素：同一组系数跨强子跨过程适用、功率修正由已知来源
主导、红外发散在 $Q$ 与 $q$ 之间相消。它还框定了后来成为标准实践的
策略——在中等 $P^z$ 提取，非微扰重整化，匹配，然后以可控系统误差外推
$\Lambda/P^z\to0$。}
\impact{与 Ji~2013 一道，这对论文定义了库内 LP3/HadStruc/ETMC 数值
论文引言所用的语言。下文的 RMP 综述是该纲领的成熟表述。}
\cardend

% ----------------------------------------------------------------------
\section{第一个匹配系数}
\papercard{Xiong, Ji, Zhang, Zhao}
{部分子分布的单圈匹配：非单态情形}
{Phys.\ Rev.\ D \textbf{90}, 014051 (2014)}
{arXiv:1310.7471 | DOI: 10.1103/PhysRevD.90.014051}
以及
\papercard{Xiangdong Ji, Jian-Hui Zhang, Yong Zhao}
{再论部分子物理的大动量有效理论方法}
{Nucl.\ Phys.\ B \textbf{924}, 366--376 (2017)}
{arXiv:1706.07416}
\whyselected{2014 年短文计算了非单态准 PDF 的首个单圈匹配核
（约 $\sim$12/99）；2017 年姊妹篇（约 $\sim$5/99）系统化了幂计数论证，
并回应了对该方法的早期质疑。两者都是"从提案到工具"不可缺少的一步。}
\physics{在维数规整下单圈分析给出显式核
\begin{equation}
  \tilde K\!\left(\frac{x}{y}\right) = \delta(x-y)
  + \frac{\alpha_s}{2\pi}\,k\!\left(\frac{x}{y}\right),
\end{equation}
$k$ 既含 AP 劈裂函数也含吸收等时与光锥运动学差异的有限项——包括通过
$\ln(P^z/\mu)$ 进入的 Wilson 线 cusp 反常量纲。NPB~924 论文证明准
算符唯一的幂发散定域于 Wilson 线自能（可因子化因而可消除），并论证
Ioffe 时间方法同样服从大 $P^z$ 逻辑，从而统一了两条提取路线。}
\impact{这些核就是把格点矩阵元变成 PDF 的东西——库内的夸克 PDF 论文
（Liu 等 2018/2020 螺旋度与非极化提取）用的正是这个非单态核。}
\cardend

% ----------------------------------------------------------------------
\section{伪 PDF 框架}
\papercard{A.\ V.\ Radyushkin}
{准部分子分布函数、动量分布与伪部分子分布函数}
{Phys.\ Rev.\ D \textbf{96}, 034025 (2017)}
{arXiv:1705.01488 | DOI: 10.1103/PhysRevD.96.034025}
\whyselected{坐标空间路线的定义性论文（约 $\sim$12/99）：库内每篇
HadStruc 论文都建立其上，ETMC 与 Lin 组也引其为共同奠基。}
\physics{Radyushkin 把同一物理重写进 Ioffe 时间空间。取
$\nu = m\sigma$（$m=2\pi/T$，格距 $\sigma$），定义 Ioffe 时间分布
\begin{equation}
  \mathcal{M}(\nu,z^2) = \langle P|\bar\psi(z)\Gamma W(z,0)\psi(0)|P\rangle,
  \qquad q(\nu)=\int_{-1}^{1}\!dx\,e^{i\nu x}q(x).
\end{equation}
格点在类空间隔 $z^2<0$ 测量 $\mathcal{M}$——即\emph{伪} ITD。除掉与
$p^2$ 形状因子同构的 $z$ 无关部分后定义\emph{约化} ITD，其不含波函数
重整化：
\begin{equation}
  \mathcal{R}(\nu,z^2) \equiv \frac{\mathcal{M}(\nu,z^2)}{\mathcal{M}(0,z^2)}
  \;\xrightarrow[|z^2|\to0]{ }\;
  \int_{-1}^1\!dx\,e^{i\nu x}q(x) + \mathcal{O}(z^2 m^2,\nu^2 m^2 z^2/\nu_0^2),
\end{equation}
这是位置空间的短距离因子化，单圈等价于 LaMET 匹配，但整个组织承袭
Balitsky--Braun 的 OPE 型推理。}
\impact{库内约化伪 ITD 提取（Fan--Lin 介子胶子、Khan 等、Egerer 等、
Salas-Chavira 等）实现的正是这个比值；Izubuchi 等人确立的两法等价性
（见下）使两条路线成为同一理论的两种方言。}
\cardend

% ----------------------------------------------------------------------
\section{首次格点实现}
\papercard{K.\ Orginos, A.\ Radyushkin, J.\ Karpie, S.\ Zafeiropoulos}
{部分子伪分布函数的格点 QCD 探索}
{Phys.\ Rev.\ D \textbf{96}, 094503 (2017)}
{arXiv:1706.05373 | DOI: 10.1103/PhysRevD.96.094503}
\whyselected{伪分布思想的首次数值实现（约 $\sim$12/99）：库内后续
HadStruc 工作与对比图共同引用的参照点。}
\physics{在淬灭及 $N_f=2$ 各向异性系综上，论文从 boosted 介子与核子
矩阵元提取约化伪 ITD，验证其对 $\nu$、$z$ 的预期标度，执行单圈匹配
进入 $x$ 空间，并在不确定度内展示与全球拟合 PDF 的一致性——确立了
可行性、误差预算（激发态、$z^2$ 截断、动量上限）与沿用至今的工作流。}
\impact{与 Radyushkin~2017 一道，本文创立了格点部分子物理的第二个
数值学派。本库中的后裔包括核子/介子/kaon 胶子伪 ITD 序列与极化探索。}
\cardend

% ----------------------------------------------------------------------
\section{因子化定理}
\papercard{T.\ Izubuchi, X.\ Ji, L.\ Jin, I.\ W.\ Stewart, Y.\ Zhao}
{联系欧氏与光锥部分子分布的因子化定理}
{Phys.\ Rev.\ D \textbf{98}, 056004 (2018)}
{arXiv:1801.03917 | DOI: 10.1103/PhysRevD.98.056004}
\whyselected{两条路线的共同严格支柱（约 $\sim$14/99）：库内论文凡需
论证"其欧氏观测量确实收敛到光锥 PDF"处皆引用之。}
\physics{利用位置空间的算符乘积展开，论文证明准 PDF 服从系数可计算的
LaMET 因子化公式，而伪 PDF 满足等价的约化 ITD 因子化；两种表述按阶经
傅里叶变换相互关联。两个对实践至关重要的结构性结果：匹配系数依赖
\emph{部分子}动量 $xP^z$ 而非仅强子动量——后来被形式化为自然标度选择
$\mu\sim\sqrt{|x|}P^z$ 并被库内的重求和论文开发利用——以及证明显式
展示了功率修正所在之处，使系统误差估计成为可能。}
\impact{此定理之后，"准 vs.\ 伪"不再是对立主张而成为交叉检验。库内
HadStruc（伪）与 LP3（准）结果的对比图正建立在这块公共地基上。}
\cardend

% ----------------------------------------------------------------------
\section{成熟的综合}
\papercard{X.\ Ji, Y.-S.\ Liu, Y.\ Liu, J.-H.\ Zhang, Y.\ Zhao}
{Large-momentum effective theory（大动量有效理论）}
{Rev.\ Mod.\ Phys.\ \textbf{93}, 035005 (2021)}
{arXiv:2004.03543 | DOI: 10.1103/RevModPhys.93.035005}
\whyselected{权威综述（约 $\sim$11/99 且持续上升）：整个框架的标准
引用，成文于领域基础已经硬化之时。}
\physics{全面综合：PDF、GPD、TMD 与分布振幅的准分布表述；完整的
重整化叙事（线性发散、辅助场、混合方案）；功率修正与 renormalons；
纵向对数重求和；截至 2020 年的数值结果现状——包括本库自己的 LP3
纲领。它还澄清了个别快报只能断言的概念要点：准观测量的普适类、
Wilson 线 cusp 的角色，以及矩方法与 $x$ 依赖方法为何互补而非对立。}
\impact{在本库内这篇综述扮演共享教科书：横率与 Boer--Mulders 论文以其
为定义性引用，智能体工作流论文采纳其流水线作为自动化的领域逻辑。
先读它会让第五、六章的每个条目更容易。}
\cardend
